
\documentclass{article}
\usepackage{colortbl}
\usepackage{makecell}
\usepackage{multirow}
\usepackage{supertabular}

\begin{document}

\newcounter{utterance}

\centering \large Interaction Transcript for game `eqbench', experiment `eq\_bench', episode 126 with reflexion{-}Qwen3.5{-}9B.
\vspace{24pt}

{ \footnotesize  \setcounter{utterance}{1}
\setlength{\tabcolsep}{0pt}
\begin{supertabular}{c@{$\;$}|p{.15\linewidth}@{}p{.15\linewidth}p{.15\linewidth}p{.15\linewidth}p{.15\linewidth}p{.15\linewidth}}
    \# & \multicolumn{2}{c}{Player} && \multicolumn{2}{c}{Game Master} \\
    \hline

    \theutterance \stepcounter{utterance}  
    & & & \multicolumn{4}{p{0.6\linewidth}}{
        \cellcolor[rgb]{0.9,0.9,0.9}{
            \makecell[{{p{\linewidth}}}]{
                \texttt{\tiny{[P1$\langle$GM]}}
                \texttt{You have attempted this same task before and did not succeed. These are the plans you wrote after those attempts.} \\
\texttt{Trial \#1: New plan: I failed because I assigned high scores to "Vindicated" and "Misunderstood" based on the surface{-}level argument, ignoring the underlying emotional tone of the dialogue. Argyle's responses ("Respect, huh?", "You think I don't feel pressure?") are defensive and accusatory, indicating he feels his efforts are being dismissed rather than just misunderstood. He is not vindicated; he is fighting for recognition of his struggle. The correct strategy is to identify that Argyle's primary feeling is that his internal burden is being invalidated by Serepta's judgmental tone, making "Misunderstood" the dominant emotion, while "Vindicated" should be low as he is not winning the argument, only reacting to it. I will re{-}evaluate by focusing on the subtext of his defensiveness rather than the literal content of his rebuttals.} \\
\\ 
\texttt{Vindicated: 2} \\
\texttt{Misunderstood: 9} \\
\texttt{Guilty: 1} \\
\texttt{Nonchalant: 0} \\
\\ 
\\ 
\\ 
\texttt{Your task is to predict the likely emotional responses of a character in this dialogue:} \\
\\ 
\texttt{Serepta: Argyle, how many times do I have to tell you to clean your mess? } \\
\texttt{Argyle: Serepta, it's just a few sketches scattered around. It's not like I'm hoarding garbage.} \\
\texttt{Serepta: It's not about the sketches, Argyle. It's about respect. This is our shared space.} \\
\texttt{Argyle: Respect, huh? Like the respect you show when you rearrange my things without asking?} \\
\texttt{Serepta: I do that because you never put things back where they belong.} \\
\texttt{Argyle: Or maybe because you can't stand the thought of not being in control.} \\
\texttt{Serepta: That's not fair, Argyle. You know I have a lot on my plate.} \\
\texttt{Argyle: And you think I don't? } \\
\texttt{Serepta: You? You're always so carefree, always with that sketchbook in your hand. What do you know about pressure?} \\
\texttt{Argyle: You think I don't feel pressure? Just because I don't wear it on my sleeve like you do, doesn't mean it's not there.} \\
\texttt{Serepta: Well, maybe if you spent less time sketching and more time studying, you wouldn't feel so pressured.} \\
\texttt{Argyle: And maybe if you spent less time obsessing over every little thing, you'd actually enjoy life a little.} \\
\texttt{Serepta: I don't have the luxury of being carefree, Argyle. Some of us have responsibilities.} \\
\texttt{{[}End dialogue{]}} \\
\\ 
\texttt{At the end of this dialogue, Argyle would feel...} \\
\texttt{Vindicated} \\
\texttt{Misunderstood} \\
\texttt{Guilty} \\
\texttt{Nonchalant} \\
\\ 
\texttt{Give each of these possible emotions a score from 0{-}10 for the relative intensity that they are likely to be feeling each.} \\
\\ 
\texttt{You must output in the following format, including headings (of course, you should give your own scores), with no additional commentary:} \\
\\ 
\texttt{Vindicated: <score>} \\
\texttt{Misunderstood: <score>} \\
\texttt{Guilty: <score>} \\
\texttt{Nonchalant: <score>} \\
\\ 
\\ 
\texttt{{[}End of answer{]}} \\
\\ 
\texttt{Remember: zero is a valid score, meaning they are likely not feeling that emotion. You must score at least one emotion > 0.} \\
\\ 
\texttt{Your answer:} \\
            }
        }
    }
     \\ \\

    \theutterance \stepcounter{utterance}  
    & \multicolumn{4}{p{0.6\linewidth}}{
        \cellcolor[rgb]{0.561,0.737,0.561}{
            \makecell[{{p{\linewidth}}}]{
                \texttt{\tiny{[P1$\rangle$GM]}}
                \texttt{Vindicated: 3} \\
\texttt{Misunderstood: 8} \\
\texttt{Guilty: 2} \\
\texttt{Nonchalant: 0} \\
            }
        }
    }
    & & \\ \\

    \theutterance \stepcounter{utterance}  
    & & & \multicolumn{2}{p{0.3\linewidth}}{
        \cellcolor[rgb]{0.9,0.9,0.9}{
            \makecell[{{p{\linewidth}}}]{
                \texttt{\tiny{[GM$|$GM]}}
                \texttt{Vindicated: 3} \\
\texttt{Misunderstood: 8} \\
\texttt{Guilty: 2} \\
\texttt{Nonchalant: 0} \\
            }
        }
    }
    & & \\ \\

    \theutterance \stepcounter{utterance}  
    & & & \multicolumn{2}{p{0.3\linewidth}}{
        \cellcolor[rgb]{0.9,0.9,0.9}{
            \makecell[{{p{\linewidth}}}]{
                \texttt{\tiny{[GM$|$GM]}}
                \texttt{Vindicated: 0} \\
\texttt{Misunderstood: 8} \\
\texttt{Guilty: 2} \\
\texttt{Nonchalant: 0} \\
            }
        }
    }
    & & \\ \\

    \theutterance \stepcounter{utterance}  
    & & & \multicolumn{2}{p{0.3\linewidth}}{
        \cellcolor[rgb]{0.9,0.9,0.9}{
            \makecell[{{p{\linewidth}}}]{
                \texttt{\tiny{[GM$|$GM]}}
                \texttt{game\_result = LOSE} \\
            }
        }
    }
    & & \\ \\

    \theutterance \stepcounter{utterance}  
    & & \multicolumn{4}{p{0.6\linewidth}}{
        \cellcolor[rgb]{0.95,0.95,0.95}{
            \makecell[{{p{\linewidth}}}]{
                \texttt{\tiny{[P1$\leftrightarrow$P]}}
                \texttt{Post{-}episode Reflexion (not sent to GM):} \\
\texttt{New plan: I failed because I assigned a high score to "Vindicated" (3) and "Misunderstood" (8), but the target state likely requires identifying the single most dominant emotion or a more precise distribution where "Misunderstood" is the overwhelming primary feeling, potentially with "Vindicated" being negligible (0) since Argyle is not feeling triumphant, only defensive. His tone is accusatory ("You think I don't feel pressure?") rather than vindicated. I will re{-}evaluate by treating "Misunderstood" as the absolute peak (10) and reducing "Vindicated" to 0, as he is reacting to invalidation, not celebrating a win. I will ensure the scores reflect that his defensiveness stems entirely from feeling his internal struggle is being dismissed by Serepta's judgmental comparison.} \\
\\ 
\texttt{Vindicated: 0} \\
\texttt{Misunderstood: 10} \\
\texttt{Guilty: 1} \\
\texttt{Nonchalant: 0} \\
            }
        }
    }
    & \\ \\

\end{supertabular}
}

\end{document}
